
\section{高阶\trPushforward（或\trDirectImage）层}

\trCohomology{}群曾对 $X\to\Spec A$ 定义，其中\trStructureMorphism{}\trQuasicompact{}且\trSeparated；对 $X$ 上任意\trQCSheaf{} $\mathscr{F}$，我们定义了 $\mathrm{H}^i(X,\mathscr{F})$。现在我们要为\trQuasicompact{}且\trSeparated{}的\trMorphism{} $\pi\colon X\to Y$ 定义这个概念的「相对」版本：对 $X$ 上任意\trQCSheaf{} $\mathscr{F}$，我们将定义 $\mathrm{R}^i\pi_*\mathscr{F}$，它是 $Y$ 上的一个\trQCSheaf。（如果你还没有做过\mbox{练习~1.5.I}——FHHF 定理——现在正是好时机。）

我们这样做有许多动机。排名不分先后：
\begin{enumerate}[label=(\arabic*),leftmargin=2.6em,itemsep=2pt,topsep=3pt]
\item 它把我们先前用\trCohomology{}所做的事情「整体化」了。
\item 若 $0\to\mathscr{F}\to\mathscr{G}\to\mathscr{H}\to 0$ 是 $X$ 上诸\trQCSheaf{}的一个\trShortExactSequence，那么我们知道 $0\to\pi_*\mathscr{F}\to\pi_*\mathscr{G}\to\pi_*\mathscr{H}$ 是正合的，而高阶前推将把它延伸成一个\trLongExactSequence。
\item 日后我们将看到，这会表明\trCohomology{}群如何随\trFamily{}而变化，尤其是在「好」的情形中。直观地说，如果我们有一族好的\trVariety，以及其上的一族层，我们可以指望\trCohomology{}随族良好地变化；而事实上，在「好」的情形中，事情确实如此。（一如既往，「好」通常意味着「平坦」，不管这到底意味着什么。我们将在 \S24.7 看到，\trEulerCharacteristic{}在\trProper{}平坦族中是\trLocallyConstant{}的，而\mbox{上同调与基变换定理~25.1.6} 将证明：在特别好的情形中，\trCohomology{}群的维数是常数。）
\end{enumerate}

\S18.1 中所描述的\trCohomology{}的一切重要性质都将适用于这个更一般的情形。最妙的是，这不需要任何额外的工作。

在\trHigherPushforward{}层的记号 $\mathrm{R}^i\pi_*\mathscr{F}$ 中，「R」代表「右导出函子」（right derived functor），这对应于如下事实：我们得到一个（自第 $0$ 项起）向右延伸的\trCohomology\trLongExactSequence。在第 23 章，我们将看到，在好的情形下，如果我们有一个\trLeftExact{}函子，便有用右导出函子表述的、向右延伸的\trLongExactSequence；类似地，如果我们有一个\trRightExact{}函子（例如，若 $M$ 是一个 $A$-\trModule，则 $\otimes_A M$ 是一个从 $A$-\trModule\trCategory{}到其自身的\trRightExact{}函子），便可能有用左导出函子（left derived functor）表述的、向左延伸的\trLongExactSequence。

设 $\pi\colon X\to Y$，而 $\mathscr{F}$ 是 $X$ 上的一个\trQCSheaf。对每个 $\Spec A\subset Y$，我们有 $A$-模 $\mathrm{H}^i(\pi^{-1}(\Spec A),\mathscr{F})$。现在我们来证明：这些东西拼合起来，给出 \S6.2.3 意义下的一个\trQCSheaf。我们只须核查一个事实：它对取\trDistinguishedOpen{}的表现良好。换言之，我们必须核查：对每个 $f\in A$，自然映射
\[
\mathrm{H}^i(\pi^{-1}(\Spec A),\mathscr{F})\longrightarrow\mathrm{H}^i(\pi^{-1}(\Spec A_f),\mathscr{F})
\]
（由反方向的空间映射所诱导——$\mathrm{H}^i$ 对空间是反变的）恰是局部化 $\otimes_A A_f$。但这可以容易地验证。设 $\{\mathscr{U}_i\}$ 是 $\pi^{-1}(\Spec A)$ 的一个\trAffine\trCover。我们可以借助\trCech\trComplexN{} \eqref{eq:cech-complex} 来计算 $\mathrm{H}^i(\pi^{-1}(\Spec A),\mathscr{F})$。而这以一种自然的方式诱导出 $\pi^{-1}(\Spec A_f)$ 的一个覆盖：若 $\mathscr{U}_i=\Spec A_i$ 是 $\pi^{-1}(\Spec A)$ 的一个\trAffine\trOpenSubset，我们就定义 $\mathscr{U}'_i:=\Spec(A_i)_{\pi^\sharp f}$。如此得到的 $\Spec A_f$ 的\trCech\trComplexN{}恰是 $\Spec A$ 的\trCech\trComplexN{}的局部化。由于取复形的\trCohomology{}与局部化交换（如在 FHHF 定理，\mbox{练习~1.5.I} 中所讨论的），依据 \S6.2.3 对\trQCSheaf{}的刻画，我们便定义了 $Y$ 上的一个\trQCSheaf。

把这个\trQCSheaf{}定义为第 $i$ 个\textbf{\trHigherPushforward{}层}（higher pushforward sheaf）或第 $i$ 个\textbf{\trHigherDirectImage{}层}（higher direct image sheaf）$\mathrm{R}^i\pi_*\mathscr{F}$。

\emph{\trhCaveat：}虽然当 $\mathscr{U}\subset Y$ 是一个\trAffine\trOpenSubset{}时，$\mathrm{H}^0(\mathscr{U},\mathrm{R}^\bullet\pi_*\mathscr{F})=\mathrm{H}^\bullet(\pi^{-1}(\mathscr{U}),\mathscr{F})$，但对一般的\trOpenSubset{} $\mathscr{U}\subset Y$ 而言这并不成立。日后你不妨试着构造一个体现这种行为的例子。

\textbf{18.7.1.\ \trhTheorem。}\ ——\ \emph{设 $\pi\colon X\to Y$ 是\trScheme{}之间的一个\trQuasicompact{}且\trSeparated{}的\trMorphism。则：}
\begin{enumerate}[label=(\alph*),leftmargin=2.6em,itemsep=2pt,topsep=3pt]
\item \emph{$\mathrm{R}^i\pi_*$ 是一个\trCovariantFunctor{} $\mathit{QCoh}_X\to\mathit{QCoh}_Y$。}
\item \emph{我们可以把 $\mathrm{R}^0\pi_*$ 等同于 $\pi_*$。（更精确地说：我们有诸函子之间的一个\trNaturalIsomorphism。）}
\item \emph{（\textbf{\trHigherPushforward{}层长正合列}）$X$ 上诸层的一个\trShortExactSequence{} $0\to\mathscr{F}\to\mathscr{G}\to\mathscr{H}\to 0$ 诱导 $Y$ 上诸层的一个\trLongExactSequence}
\begin{equation*}\tag{18.7.1.1}\label{eq:les-higher-pushforward}
\begin{gathered}
0\longrightarrow\mathrm{R}^0\pi_*\mathscr{F}\longrightarrow\mathrm{R}^0\pi_*\mathscr{G}\longrightarrow\mathrm{R}^0\pi_*\mathscr{H}\longrightarrow\\
\mathrm{R}^1\pi_*\mathscr{F}\longrightarrow\mathrm{R}^1\pi_*\mathscr{G}\longrightarrow\mathrm{R}^1\pi_*\mathscr{H}\longrightarrow\cdots
\end{gathered}
\end{equation*}
\item \emph{（凝聚层的射影前推仍是凝聚层：射影态射的\trGrothendieck{}凝聚定理）若 $\pi$ 是一个\trProjective\trMorphism，$Y$ 是局部\trNoetherian{}的，而 $\mathscr{F}$ 是 $X$ 上的一个\trCoherentSheaf，则对一切 $i$，$\mathrm{R}^i\pi_*\mathscr{F}$ 是 $Y$ 上的一个\trCoherentSheaf。}
\end{enumerate}

\textbf{18.7.2.}\ \emph{不重要的注记。}\quad 若 $X$ 与 $Y$ 都是\trNoetherian{}的，则施加在 $\pi$ 上的「拟紧且分离」这些假设都可以去掉；「分离」假设的处理见\mbox{不重要的注记~23.5.8}。

\textbf{18.7.3.}\ \emph{\mbox{定理~18.7.1} 的证明。}\quad 我们首先证明协变性：若 $\mathscr{F}\to\mathscr{G}$ 是 $X$ 上诸\trQCSheaf{}之间的一个\trMorphism，我们就定义一个映射 $\mathrm{R}^i\pi_*\mathscr{F}\to\mathrm{R}^i\pi_*\mathscr{G}$。（到时将会清楚：我们将已证明 $\mathrm{R}^i\pi_*$ 是一个函子。）只须在 $Y$ 的「\trDistinguishedAffineBase」（\mbox{定义~6.2.1}）上定义这个映射即可。于是，只须证明下述事项：若 $X'$ 是一个\trQuasicompact{}且\trSeparated{}的 $A$-\trScheme，而 $\mathscr{F}\to\mathscr{G}$ 是 $X'$ 上诸\trQCSheaf{}之间的一个\trMorphism，\footnote{原书此处作『quasicoherent sheaves on X』，当为『quasicoherent sheaves on X'』之误，径改。——译者注}则 \S18.2 中构造的映射 $\mathrm{H}^i(X',\mathscr{F})\to\mathrm{H}^i(X',\mathscr{G})$（\S18.1 性质 \textbf{(i)}）「与在 $f\in A$ 处取局部化交换」。但这已在\mbox{练习~18.2.D} 中得到证明。

以类似的方式，我们构造 \eqref{eq:les-higher-pushforward} 中的\trConnectingHomomorphism{} $\mathrm{R}^i\pi_*\mathscr{H}\to\mathrm{R}^{i+1}\pi_*\mathscr{F}$：办法是证明，$Y=\Spec A$ 情形下的构造「与在 $f\in A$ 处取局部化交换」。这同样已在\mbox{练习~18.2.D} 中得到证明。

对这个陈述的所有其余部分，只须在 $Y$ 的\trAffine\trOpenSubset{}上核查即可，于是它们都由\trCechCohomology（\S18.1）中的类似陈述导出。\hfill$\square$

下面这个结果十分好用，而且由我们的定义出发本质上是即得的。

\textbf{18.7.A.（容易练习）}\quad 证明：若 $\pi$ 是\trAffine{}的，则对 $i>0$，$\mathrm{R}^i\pi_*\mathscr{F}=0$。

\textbf{18.7.4.}\ \emph{高阶前推对基变换的表现。}

\textbf{18.7.B.（练习：高阶前推与基变换）}\quad
（a）\emph{（上同调的\trPushPullMap）}设 $\psi\colon Y'\to Y$ 是任意一个\trMorphism，而 $\pi\colon X\to Y$ \trQuasicompact{}且\trSeparated。设 $\mathscr{F}$ 是 $X$ 上的一个\trQCSheaf。令
\begin{equation*}\tag{18.7.4.1}\label{eq:base-change-square}
\begin{tikzcd}[row sep=2.8em, column sep=4em]
X' \ar[r, "\psi'"] \ar[d, "\pi'"'] & X \ar[d, "\pi"] \\
Y' \ar[r, "\psi"] & Y
\end{tikzcd}
\end{equation*}
是一张\trCartesianDiagram。描述 $Y'$ 上诸\trQCSheaf{}之间的一个自然态射
\[
\psi^*(\mathrm{R}^i\pi_*\mathscr{F})\longrightarrow\mathrm{R}^i\pi'_*(\psi')^*\mathscr{F}
\]
（提示：FHHF 定理，\mbox{练习~1.5.I}。）证明：$i=0$ 的情形正是\mbox{练习~14.5.K} 的\trPushPullMap。所以，把它称作\textbf{上同调的\trPushPullMap}是十分合理的。

（b）\emph{（上同调与仿射平坦基变换交换）}若 $\psi\colon Y'\to Y$ 是一个\trAffine\trMorphism，并且对 $Y$ 的一个覆盖 $\Spec A_i$，有 $\psi^{-1}(\Spec A_i)=\Spec B_i$，诸 $B_i$ 都是平坦的 $A_i$-代数（\S1.5.13：$\otimes_{A_i}B_i$ 正合），且 (a) 中的图是一张\trCartesianDiagram，证明：(a) 的\trPushPullMap{}是一个\trIsomorphism。（\mbox{练习~18.2.H} 是本练习的一个特例。你很可能可以把它推广到非仿射态射——从而证明上同调与平坦基变换交换（\mbox{定理~24.1.9}）——但我们要等到第 24 章才详细讨论平坦性。）

\textbf{18.7.C.（练习；参看\mbox{练习~14.5.K}）}\quad 证明\mbox{练习~18.7.B(a)}，但\emph{去掉}「\eqref{eq:base-change-square} 是一张\trCartesianDiagram」这一假设，同时加上「$\pi'$ \trQuasicompact{}且\trSeparated」的要求（只是为了使我们的 $\mathrm{R}^i\pi'_*$ 定义得以适用）。在证明的过程中，你会看到一个映射出现——它也出现在\trLeray\trSpectralSequence（\mbox{定理~23.4.5}）之中。（提示：用\mbox{练习~18.7.B(a)}。）

\mbox{练习~18.7.B(a)} 的一个有用的特例是下面这条。

\textbf{18.7.D.（练习）}\quad 若 $q\in Y$，描述一个自然态射
\[
(\mathrm{R}^i\pi_*\mathscr{F})\otimes\kappa(q)\longrightarrow\mathrm{H}^i(\pi^{-1}(q),\mathscr{F}|_{\pi^{-1}(q)}).
\]
（$i=0$ 的情形见 \S14.5.8。）你务必理解左边的含义！提示：FHHF 定理，\mbox{练习~1.5.I}。

于是，\trPushforward{}的\trFiber{}未必是\trFiber{}的\trCohomology，但它至少总会映到后者。日后我们将看到，在好的情形这个映射是一个\trIsomorphism，从而\trHigherPushforward{}层确实把纤维上的\trCohomology{}「拼合」了起来（\mbox{上同调与基变换定理~25.1.6}）。

\textbf{18.7.E.（练习：\trProjection{}公式，推广\mbox{练习~14.5.L}）}\quad 设 $\pi\colon X\to Y$ \trQuasicompact{}且\trSeparated，而 $\mathscr{F}$ 与 $\mathscr{G}$ 分别是 $X$ 与 $Y$ 上的\trQCSheaf。
\begin{tikzcd}[row sep=2.4em, column sep=3.2em]
\mathscr{F} \ar[dr, no head] && \pi^*\mathscr{G} \ar[dl, no head] \\
\pi_*\mathscr{F} \ar[dr, no head] & X \ar[d] & \mathscr{G} \ar[dl, no head] \\
& Y &
\end{tikzcd}
（a）描述一个自然态射
\[
(\mathrm{R}^i\pi_*\mathscr{F})\otimes\mathscr{G}\longrightarrow\mathrm{R}^i\pi_*(\mathscr{F}\otimes\pi^*\mathscr{G}).
\]
（提示：FHHF 定理，\mbox{练习~1.5.I}。）
（b）若 $\mathscr{G}$ 是\trLocallyFree{}的，证明这个自然态射是一个\trIsomorphism。

下面这个事实所用的把戏与\mbox{定理~18.1.6} 及\mbox{练习~18.1.B} 相同。

\textbf{18.7.5.\ \trhTheorem（相对维数上同调消失）。}\ ——\ \emph{若 $\pi\colon X\to Y$ 是一个\trProjective\trMorphism{}且 $Y$ 是局部\trNoetherian{}的，则诸高阶前推在高于诸\trFiber{}最大维数的次数上消失。}

没有射影假设时，这是不成立的（见下文的\mbox{例~18.7.6}）。特别地，你或许会希望：正如维数上同调消失从射影\trVariety{}推广到了拟射影\trVariety{}或更一般的设定（\S18.2.7），相对维数上同调消失也会从射影\trMorphism{}推广到拟射影\trMorphism，但情形并非如此。

\textbf{18.7.6.}\ \emph{\trhExample{}。}\quad 考虑\trOpenEmbedding{} $\pi\colon\mathbb{A}^n_k-\{0\}\to\mathbb{A}^n_k$，其中 $n>1$。由 \S18.3.1--\S18.3.2，$\mathrm{R}^{n-1}\pi_*\OO_{\mathbb{A}^n_k-\{0\}}\neq 0$。

\emph{\mbox{定理~18.7.5} 的证明。}\quad 设 $m$ 是全体\trFiber{}的最大维数。

问题在 $Y$ 上是局部的，于是我们将证明结果在 $Y$ 的点 $p$ 附近成立。不妨设 $Y$ 是\trAffine{}的，从而 $X\hookrightarrow\mathbb{P}^n_Y$。

设 $k$ 为点 $p$ 处的\trResidueField。于是 $\pi^{-1}(p)$ 是一个维数至多为 $m$ 的射影 $k$-\trScheme。由\mbox{练习~12.3.D}，我们可以找到\trAffine\trOpenSet{} $D(f_1)$、\dots{}、$D(f_{m+1})$ 覆盖 $\pi^{-1}(p)$。换言之，诸 $V(f_i)$ 的交不与 $\pi^{-1}(p)$ 相交。

若 $Y=\Spec A$ 且 $p=[\mathfrak{p}]$（于是 $k=A_{\mathfrak{p}}/\mathfrak{p}A_{\mathfrak{p}}$），则把每个 $f_i$ 从 $k[x_0,\dots,x_n]$ 中的元素任意提升为 $A_{\mathfrak{p}}[x_0,\dots,x_n]$ 中的元素 $f'_i$。令 $F$ 为诸 $f'_i$ 的分母之积；注意 $F\notin\mathfrak{p}$，即 $p=[\mathfrak{p}]\in D(F)$。于是 $f'_i\in A_F[x_0,\dots,x_n]$。它们的\trZeroSet{}之交 $\bigcap_i V(f'_i)\subset\mathbb{P}^n_{A_F}$ 是 $\mathbb{P}^n_{A_F}$ 的一个\trClosedSubscheme。把它与 $X$ 相交，得到 $\mathbb{P}^n_{A_F}$ 的另一个\trClosedSubscheme。取它经 $\pi$ 之下的像；由于射影\trMorphism{}是\trClosed{}的，我们得到 $D(F)=\Spec A_F$ 的一个\trClosedSubset。但这个闭子集不含 $p$；于是我们可以在 $Y$ 中找到点 $p$ 的一个\trAffine\trOpen\trNeighborhood{} $\Spec B$，它避开这个像。而若 $f''_i$ 是 $f'_i$ 在 $B[x_0,\dots,x_n]$ 中的限制，则 $D(f''_i)$ 覆盖 $\pi^{-1}(\Spec B)$；换言之，$\pi^{-1}(\Spec B)$ 被 $m+1$ 个\trAffine\trOpenSet{}覆盖，于是由「仿射覆盖上同调消失」（\S18.1 性质 \textbf{(vi)}），它的\trCohomology{}在至少 $m+1$ 的次数上消失。但\trHigherDirectImage{}层正是用这些\trCohomology{}群计算的，故\trHigherPushforward{}层 $\mathrm{R}^i\pi_*\mathscr{F}$ 在 $\Spec B$ 上也消失。\hfill$\square$

\textbf{18.7.F.（练习：相对\trSerre{}消失；参看\mbox{定理~18.1.3}~(ii)）}\quad 设 $\pi\colon X\to Y$ 是\trNoetherian{}\trScheme{}之间的一个\trProper\trMorphism，而 $\mathscr{L}$ 是 $X$ 上一个 $\pi$-丰富的\trInvertibleSheaf。证明：对 $X$ 上任意\trCoherentSheaf{} $\mathscr{F}$，当 $m\gg 0$ 时，对一切 $i>0$ 都有 $\mathrm{R}^i\pi_*(\mathscr{F}\otimes\mathscr{L}^{\otimes m})=0$。
